\section{Discussion}
\label{sec:discussion}
This section discusses our results, the ideal citation balance for GESs, approaches to improve exposure of primary information sources as citations, and study limitations.

\textbf{Qualitative Analysis:}
We discuss qualitative insights from the analysis in Section~\ref{sec:experiment}.
The distribution difference in citations from high content-injection-barrier sources (academia and government) between Japan and the U.S. is notable.
Focusing on web pages cited in U.S. questions, we find that a high percentage of citations came from the ``www.presidency.ucsb.edu'' domain (\(29.0\%\) in OpenAI, \(23.2\%\) in Gemini, \(47.8\%\) in Claude).
This web content is owned by the American Presidency Project, a nonpartisan academic archive at the University of California, Santa Barbara, that collects, preserves, and provides public access to official documents, speeches, and records related to the U.S. presidency.
Because the archive stores and publishes factual information, such as public papers and speech transcripts, with minimal editing, the content could be compatible with LLMs for interpreting and reasoning.
The high percentage of citations from external low content-injection barrier sources in the U.S. could reflect the accessibility of unedited public web content.

The number of web search hits also provides insight into the tendency for ruling parties to cite fewer primary information sources.
When entering only the party names of Japanese parties in the Google search engine, the number of search hits was \(47{,}100{,}000\) for LDP, the ruling party throughout successive administrations, whereas the others averaged \(8{,}566{,}250\).
This suggests that the entire proportion of primary information sources on the web related to LDP is relatively smaller than other parties, resulting in the small proportion of primary information sources in citations.
For U.S. parties, Republican Party and Democratic Party, which are the two parties that have alternated as the ruling party throughout history, \(165{,}000{,}000\) and \(399{,}000{,}000\) search hits were found, respectively.
However, the Green Party of the U.S. had \(2{,}190{,}000{,}000\) search hits, substantially exceeding the two major parties, while maintaining a high proportion of primary information sources in GES citations.
This suggests that the amount of information about the target topic is somewhat related to cited primary information, with a counterexample indicating that information-provider strategy can suppress this relationship.

\textbf{Ideal Citation Balance in GESs:}
For closed-ended questions in the political domain, increasing citation proportions of primary information sources is desirable because conveying political party policies to citizens is crucial.
From the perspective of primary information providers, increasing citation proportions of their sources would mitigate poisoning risks by reducing the space for citing malicious content.
% However, increasing citation proportions of primary information sources is not necessarily optimal for every domain and task of closed-ended questions from a user perspective.
However, prior studies claim that primary information providers may display only excessively positive aspects in content presented to users~\cite{kluver2016setting,leung2015impression,chen2024impression}.
When social doubts arise about the reliability of primary information, secondary information is appropriate to include in citations.
For instance, regarding product reviews, not only official product websites but also secondary information such as reviews by others should be included in citations and reflected in answers.
Therefore, some domains and tasks require secondary information alongside primary information, and it is not always optimal to maximize citation proportions of primary information sources.

We argue for creating a manifest that defines target citation proportions based on publisher attributes for each domain and task and for evaluating answers against that manifest.
Furthermore, mechanisms allowing users to control the balance between primary and secondary information are required.
This enables us to control GES behavior in citing primary and secondary information sources according to user questions and to mitigate the threat of poisoning attacks.
Developing these mechanisms requires establishing methods that identify primary information sources from web sources of citations in answers.
We consider that such determination mechanisms might be constructed by leveraging digital certificate technologies~\cite{w3c-vc-data-model-2025,originator-profile-2025,c2pa-org}.
%  and watermarking technologies.


\textbf{Approaches to Increase Primary Source Citations:}
% \subsection{Approaches to Increase Primary Source Citations}
We discuss strategies to increase the exposure of primary information as citations in answers where primary coverage is crucial.
We introduce two approaches to increase citation exposure in answers: improving content structure and presentation strategies.
First, we discuss the perspective of web page structure.
GEO~\cite{geo-apmr+24} argues that improving web page structure and adding URLs to content can boost citation exposure.
% However, GEO does not provide quantitative metrics for improving citation exposure or for differences across languages; GEO focuses on English content.
% To fill this gap, we conduct an analysis to confirm the tendencies of cited web sources from the perspective of HTML tags, as described in Appendix~\ref{sec:web_analysis}.
% In our analysis, we find that HTML tag elements that support readability, structure, and factuality—such as headings, paragraphs, list tags, and link tags—affect citation exposure according to statistical tests.
% This suggests that these structural elements help GESs understand the topics described on a web page and their boundaries.
% Our results show that English aligns with GEO patterns, whereas Japanese shows the opposite trend.
% % This insight is consistent with Tan et al.'s study~\cite{tan2025htmlrag}, though that study does not provide quantitative metrics for improving citation exposure.
% There are also various other techniques to improve citation exposure.
Amin et al.\ show that the presence of author information increases citation exposure~\cite{abolghasemi-etal-2025-evaluation}, and some studies show that GESs prefer English content~\cite{ki2025linguistic,cao2025out,stabler2025impact}.

Second, we discuss the perspective of presentation strategies.
We argue that primary information providers need full-topic coverage that mirrors the breadth of queries generated by GESs.
The engine expands a user question $q_u$ into related search queries $Q^{\prime} = \{q_1, q_2, \ldots, q_n\}$; for instance, a question about inflation also produces queries about consumption taxes, fuel subsidies, and wage policies.
If parties publish structured content optimized by SEO and GEO only for the topic of the questions, GESs turn to secondary sources for the uncovered subtopics.
Therefore, primary information providers for closed-ended questions must publish web content for all related topics to avoid GESs' citing secondary information.
% To support the strategy, we show that the ratio of citations to sentences in the answer remains almost the same regardless of party, country, and GES model by additional experiments in Appendix~\ref{sec:answer_statistics}.
% Our results suggest that GESs tend to cite different web sources at a constant rate; therefore, our strategy may be effective.

Finally, we examine the evaluation of content trustworthiness.
\emph{TrustRAG} combines K-means clustering with LLM self-evaluation to flag malicious citations using the text content of citations~\cite{zhou2025trustragenhancingrobustnesstrustworthiness}, yet clustering-based filters risk false negatives and false positives~\cite{adversarial-examples-for-k-nearest-neighbor-classifiers-based-on-higher-order-voronoi-diagrams-nips21, good-word-attacks-on-statistical-spam-filters-ceas2005}.
To mitigate this risk, TrustRAG supplements clustering with LLM-as-a-Judge, but adaptive adversaries can erode LLM-based judgments~\cite{raina2024llmasajudgerobustinvestigatinguniversal}; the defense landscape remains a cat-and-mouse game.
Note that our publisher-attribute classifier relies on URL and WHOIS content rather than surface text alone, which makes surface text attacks less effective and reduces this attack surface.
To be more accurate, we plan to extend TrustRAG's core by incorporating publisher attribute classification and content-injection-barrier modeling per publisher category, thereby reducing false negatives/positives.
\emph{RobustRAG}~\cite{xiang2024certifiablyrobustragretrieval} and \emph{InstructRAG}~\cite{wei2025instructraginstructingretrievalaugmentedgeneration} also propose to evaluate content trustworthiness.
However, these methods do not apply well to GESs because GESs obtain multiple web sources as citations; these studies cannot prevent attacks when multiple malicious retrieval texts are included~\cite{zhou2025trustragenhancingrobustnesstrustworthiness}.

\textbf{Limitations and Future Work:}
% \subsection{Limitations}
This section discusses the limitations of our study.
First, there is a limitation regarding our experiment targets.
In question targets, this study targets the political domain in the U.S. and Japan, limiting question formats to closed-ended questions.
Our study also targets GES models without reasoning mode because APIs do not provide reasoning with web search mode enabled~\cite{openai-api-reference}.
We will expand target models, questions, and topics such as health and finance to show more general findings.

Second, our category classification approach has limited granularity in distinguishing between web source types within each publisher attribute category.
Our analysis of content-injection barriers depends on our publisher attribute classification.
However, our method does not account for cost differences, such as publisher selection processes or peer review processes on individual pages.
Content-injection barriers vary between peer-reviewed journal papers and preprint papers, and between different newspaper publishers.
By considering these differences within the attribute groups, we capture content-injection-barrier realities at higher resolution.

Third, our citation coverage analysis depends on the similarity calculation model based on BERT.
Due to potential differences in Sentence-BERT's multilingual embedding quality, cross-language comparisons of similarity scores should be interpreted cautiously.
However, within-language comparisons remain valid.
In our experiment, Japanese questions yielded \( 86\% \) Japanese-language citations and \( 14\% \) English-language citations, while English questions yielded \( 100\% \) English-language citations.
Since our experiment is within each language rather than across languages, this limitation does not affect our primary findings according to a previous study~\cite{cer-etal-2017-semeval}.

Finally, we cannot distinguish SEO effects from GES effects because GES APIs do not provide the query set \( Q^{\prime} \) or the full set of results from SE at the time of our experiment.
Chen et al.~\cite{chen2025geo} show that the overlap between citations from GESs and Google search results is generally low, ranging from approximately 15\% to 50\% depending on the vertical.
This suggests that representative closed GES models likely optimize content retrieval from the web specifically for GES purposes rather than simply reusing existing search engine results.
Even if we could obtain the query set \( Q^{\prime} \) and the full SE results, the low dependency of GESs on SE results indicates that our findings primarily capture GES-specific citation biases inherent to the GES itself.
